% ============================================================================
% II. RELATED WORK (REVISED)
% ============================================================================
\section{Related Work}
\label{sec:related_work}

This section reviews prior work in four areas relevant to the present framework: hyperspectral imaging systems and data characteristics, band selection methods for hyperspectral data, deep learning approaches for spectral analysis, and explainability methods for feature attribution.

% ----------------------------------------------------------------------------
\subsection{Multi-Excitation Hyperspectral Imaging Systems}
\label{sec:related_hsi}

Hyperspectral imaging systems have evolved from early pushbroom scanners requiring platform motion~\cite{goetz1985imaging} to contemporary snapshot sensors capable of real-time acquisition~\cite{hagen2013review}.
While reflectance-based hyperspectral imaging remains dominant for remote sensing applications, fluorescence hyperspectral imaging has gained traction in laboratory and close-range settings where controlled excitation is feasible~\cite{mishra2017close}.
The extension to multiple excitation wavelengths---yielding the EEM structure that characterizes ME-HSI---was pioneered in analytical chemistry for fluorescence spectroscopy of solutions~\cite{murphy2013fluorescence} and has been adapted for imaging applications where spatially resolved EEM data provides additional discrimination power~\cite{lavine2001exploratory}.
Such systems have proven particularly valuable for biological specimens, including lichens and plant tissues, where autofluorescence signatures encode information about molecular composition that reflectance imaging cannot capture~\cite{croce2019autofluorescence}.

The 4D data structure of ME-HSI presents unique challenges.
Unlike 3D hyperspectral cubes where spectral bands sample a single physical phenomenon (reflectance or emission at fixed excitation), ME-HSI data exhibits complex cross-excitation correlations arising from energy transfer between fluorophores, concentration-dependent quenching, and overlapping emission spectra~\cite{lakowicz2006principles}.
Additionally, Rayleigh and Raman scattering at wavelengths near the excitation produce artifacts that must be removed, resulting in variable numbers of valid emission bands per excitation~\cite{murphy2013fluorescence}.
While hardware advances have made 4D acquisition practical, analytical methods for handling the resulting data dimensionality have lagged behind~\cite{mishra2017close}.

% ----------------------------------------------------------------------------
\subsection{Band Selection for Hyperspectral Data}
\label{sec:related_band}

Band selection methods aim to identify a subset of spectral bands that retain discriminative information while reducing dimensionality~\cite{sun2019hyperspectral}.
These methods are typically categorized as filter, wrapper, or embedded approaches.

Filter methods rank bands based on statistical properties computed without reference to any downstream task.
Variance-based selection prioritizes bands with high variance across samples~\cite{chang2000constrained}, entropy-based methods favor bands with maximum information content~\cite{guo2006band}, and correlation-based approaches remove redundant bands that are highly correlated with others~\cite{martinez2006comparison}.
While computationally efficient, filter methods lack any notion of discriminative value---a band with high variance may capture noise rather than meaningful variation.
Moreover, they treat bands independently, ignoring the joint information captured by band combinations.

Wrapper methods evaluate band subsets based on the performance of a classifier trained on the selected features.
Sequential forward selection (SFS) greedily adds bands that most improve classification accuracy~\cite{pudil1994floating}, while genetic algorithms and particle swarm optimization explore the combinatorial search space more globally~\cite{pal2010feature}.
These methods can identify task-relevant bands but suffer from computational cost that grows combinatorially with the number of bands.
For a 192-band ME-HSI dataset, exhaustive evaluation of all possible subsets is intractable.

Embedded methods integrate band selection into model training.
Sparse regression (LASSO) induces feature selection through $L_1$ regularization~\cite{tibshirani1996regression}, attention mechanisms in neural networks learn to weight input bands~\cite{mou2017deep}, and reinforcement learning formulations treat band selection as a sequential decision process~\cite{roy2020hybridsn}.
While elegant, these approaches typically require labeled training data, limiting their applicability when ground truth is scarce or expensive to obtain.

Critically, nearly all existing band selection methods assume a 3D data structure where bands can be evaluated independently.
For ME-HSI, where information lies in excitation-emission \emph{pairs} and cross-excitation patterns may be discriminative, methods that explicitly model the 4D structure are needed.
Recent work on tensor decomposition for EEM data~\cite{murphy2013fluorescence} and multi-way analysis~\cite{bro1997parafac} addresses spectroscopic rather than imaging data.
The present framework bridges this gap by extending deep learning band selection to the 4D imaging case.

% ----------------------------------------------------------------------------
\subsection{Deep Learning for Hyperspectral Analysis}
\label{sec:related_dl}

Deep learning has substantially advanced hyperspectral image analysis, with convolutional neural networks achieving state-of-the-art results in classification, segmentation, and anomaly detection~\cite{li2019deep,paoletti2019deep}.
Early approaches applied 2D CNNs to spectral vectors treated as 1D signals~\cite{hu2015deep}, while subsequent work exploited spatial context through 3D convolutions that jointly process spatial and spectral dimensions~\cite{chen2016deep}.
Architectures such as spectral-spatial residual networks~\cite{zhong2017spectral} and HybridSN~\cite{roy2020hybridsn} combine 3D convolutions for spectral-spatial feature extraction with 2D convolutions for spatial refinement, achieving classification accuracies exceeding 95\% on standard benchmarks.

For unsupervised analysis, autoencoders have been applied to hyperspectral data for compression~\cite{chen2014deep}, denoising~\cite{vincent2008extracting}, and feature extraction~\cite{mou2017unsupervised}.
Variational autoencoders (VAEs) provide probabilistic latent representations that enable uncertainty quantification~\cite{kingma2013auto}, while sparse autoencoders encourage disentangled features through sparsity constraints~\cite{ng2011sparse}.
Stacked convolutional autoencoders~\cite{masci2011stacked} have demonstrated the ability to learn hierarchical features from spectral data without supervision.
These works focus on learning representations for downstream classification; in contrast, the present framework uses the learned representation as the basis for wavelength attribution.

Most deep learning approaches for hyperspectral data assume a 3D input structure.
While 3D convolutions naturally extend to 4D data by adding an excitation dimension, the variable number of emission bands per excitation after Rayleigh cutoff complicates a monolithic architecture.
The parallel-branch design introduced here addresses this limitation, enabling end-to-end learning on 4D ME-HSI data with heterogeneous spectral dimensions.

% ----------------------------------------------------------------------------
\subsection{Explainability and Attribution Methods}
\label{sec:related_explain}

Explaining neural network predictions has become crucial for scientific applications where understanding \emph{why} a model makes decisions is as important as accuracy~\cite{gilpin2018explaining}.
Post-hoc explanation methods include saliency maps that compute gradients of outputs with respect to inputs~\cite{simonyan2014deep}, Grad-CAM that localizes class-discriminative regions~\cite{selvaraju2017grad}, LIME that fits local linear models~\cite{ribeiro2016should}, and SHAP that provides game-theoretic feature attributions~\cite{lundberg2017unified}.

Perturbation-based methods offer a model-agnostic approach to attribution: by perturbing inputs or intermediate representations and observing changes in output, one can infer which features are important without access to model internals~\cite{zeiler2014visualizing}.
This approach has been applied to image classification by occluding image regions~\cite{zeiler2014visualizing}, to text analysis by removing words~\cite{li2016understanding}, and to tabular data by shuffling feature values~\cite{fisher2019all}.

In the context of autoencoders, latent space perturbation provides insights into what the model has learned~\cite{higgins2017beta}.
By perturbing individual latent dimensions and observing reconstruction changes, one can identify which input features are encoded by each dimension.
This principle underlies disentanglement analysis in VAEs~\cite{chen2018isolating} and has been used for concept discovery in generative models~\cite{voynov2020unsupervised}.
The present work adapts latent perturbation for wavelength selection in ME-HSI: rather than seeking disentangled concepts, reconstruction sensitivity is traced back to specific excitation-emission combinations to identify informative wavelengths.

% ----------------------------------------------------------------------------
\subsection{Summary and Positioning}
\label{sec:related_summary}

In summary, existing band selection methods predominantly assume a 3D hyperspectral data structure and rely on either statistical properties (filter methods), supervised classification performance (wrapper methods), or labeled training data (embedded methods).
Deep learning approaches for hyperspectral imaging have focused primarily on supervised classification rather than unsupervised wavelength analysis.
Meanwhile, explainability methods have been developed for understanding model predictions but have not been systematically applied to spectral band selection.

The proposed framework addresses these gaps by: (1)~designing a 4D-aware autoencoder architecture with parallel excitation branches that accommodates variable band counts, (2)~leveraging unsupervised representation learning to avoid dependence on class labels, and (3)~applying perturbation-based attribution to trace learned importance back to interpretable wavelength selections.
This combination enables principled wavelength selection for ME-HSI data without requiring ground-truth annotations.